
%
%Modelo com instruções para preparação de trabalho para o VI EIC 2020
%

\documentclass[12pt]{article}
\usepackage[utf8]{inputenc}
\usepackage[brazil]{babel}
\usepackage{graphicx}
\usepackage{float}
\usepackage{epsfig}
\usepackage{mathptmx} %para Times Roman
\usepackage{setspace} \setstretch{1.15} %para espaçamento 1.15
\usepackage{geometry} %para definir a geometria da página
\usepackage{booktabs} %pacote para tabelas
\usepackage{caption} %Legendas das figuras
\geometry{left=3cm,right=2cm, top=7cm, bottom=2cm} %margens da página
\setlength{\parindent}{1cm} %tamanho do parágrafo
\usepackage{indentfirst} %primeiro parágrafo do \section 
\usepackage{background}
\backgroundsetup{contents={\includegraphics[scale=1.4]{logo}},
	scale=1,placement=top,opacity=0.5,position={8.3cm,1cm}, vshift=35}
\pagestyle{empty}%retirando numeração


\everymath{\displaystyle}
\newtheorem{exem}{Exemplo}

\begin{document}
	
\begin{center}
\textbf{\Large TÍTULO (CAIXA ALTA): subtítulo (se houver) de forma clara, objetiva e em caixa baixa.}
\end{center}	

\begin{center}
\textit{	
	Nome do Autor\footnote{Biografia suscinta, E-mail de contato. (Informação Opcional).}, 
	Nome do Autor\footnote{Biografia suscinta, E-mail de contato. (Informação Opcional).}, 
	Nome do Autor\footnote{Biografia suscinta, E-mail de contato. (Informação Opcional).}, \\
	Nome do Autor  (Orientador)\footnote{Biografia suscinta, E-mail de contato. (Informação Opcional).}, 
	Nome do Autor  (Coorientador)\footnote{Biografia suscinta, E-mail de contato. (Informação Opcional).}
}
\end{center}

\begin{center}
	Nome da Instituição ou Escola, Cidade, UF
\end{center}


\begin{center}
\textit{	Eixo Temático  nº $(I - V)$: Nome do Eixo\\
		Desafio nº $(1 - 3)$: Nome do desafio  / (Edital 2) Resumo expandido}
\end{center}


{\setlength{\parindent}{0cm} %tamanho do parágrafo


\begin{center}
	\textbf{RESUMO}
\end{center}



O resumo deve ser elaborado conforme a ABNT NBR 6028:2003. Resumo é uma apresentação concisa dos pontos relevantes de um documento. O resumo deve ressaltar o objetivo, o método, os resultados e as conclusões do documento. A ordem e a extensão destes itens dependem do tipo de resumo (informativo ou indicativo) e do tratamento que cada item recebe no documento original. O resumo deve ser composto de uma sequência de frases concisas, afirmativas e não de enumeração de tópicos. Recomenda-se o uso de parágrafo único. A primeira frase deve ser significativa, explicando o tema principal do documento. A seguir, deve-se indicar a informação sobre a categoria do tratamento (memória, estudo de caso, análise da situação etc.). Deve-se usar o verbo na voz ativa e na terceira pessoa do singular. 



\vspace{2ex}

\textbf{Palavras-chave}: Palavra 1. Palavra 2. Palavra 3. Palavra 4. Palavra 5.




\begin{center}
	\textbf{ABSTRACT}
\end{center}


The abstract must be prepared in accordance with ABNT NBR 6028:2003. Summary is a concise presentation of the relevant points of a document. The summary should highlight the purpose, method, results and conclusions of the document. The order and extent of these items depend on the type of summary (informative or indicative) and the treatment each item receives in the original document. The abstract should be composed of a sequence of concise, affirmative and nontopic phrases. It is recommended to use single paragraph. The first sentence should be meaningful, explaining the main theme of the document. The information on the treatment category (memory, case study, situation analysis, etc.) should be indicated in sequence. The active voice and verbs in the third person singular must be used.

\vspace{2ex}

\textbf{Keywords:} Word 1. Word 2. Word 3. Word 4. Word 5.

}





\section{\large INTRODUÇÃO}


Na Parte introdutória do texto se deve dizer ao leitor o que ele pode esperar do conjunto do trabalho. A introdução deve fazer referência\footnote{Para inserir notas de rodapé deve-se seguir as orientações da norma ABNT NBR 10520:2002.}  ao contexto em que o problema da pesquisa está inserido, por exemplo, se é uma demanda social, um desafio tecnológico, uma curiosidade científica. Depois deve apontar com clareza o problema central e os problemas secundários que foram pesquisados.

Por fim apontar as condições em que o projeto foi desenvolvido, se ele se limitou a escola, se houve procura de informações em outros lugares, se teve apoio de outras pessoas ou instituições. Introdução é um texto corrido contendo uma explanação sobre o tema que você pesquisou ou interviu, pautado nas ideias dos autores estudados durante a pesquisa e também as informações extraídas do seu plano de pesquisa como problemática, objetivos e hipótese. 

O texto submetido deverá conter o mínimo de 3 (três) e o máximo de 5 (cinco) laudas contando com as referências. 

As citações (conforme a ABNT NBR 10520:2002) com mais de três linhas, paginação, notas, legendas e fontes das ilustrações e tabelas devem ser em tamanho menor e uniforme. Os títulos das seções com ou sem indicativo numérico devem ser conforme a ABNT NBR 6024. 

A sigla, quando mencionada pela primeira vez no texto, deve ser indicada entre parênteses, precedida do nome completo, exemplo: Associação Brasileira de Normas Técnicas (ABNT).


\section{\large MATERIAIS E MÉDOTOS}

Neste ponto deve ser feita a descrição das fontes de consultas e as metodologias utilizadas na realização da pesquisa. Ou seja, se houve pesquisa bibliográfica em bibliotecas, experimentação em laboratórios, consulta a websites, entrevistas ou outras formas de obter as informações necessárias ao desenvolvimento da pesquisa.

Descreva como o trabalho foi realizado (procedimentos, estratégias; os sujeitos, participantes, documentos; equipamentos, ambientes; etc.).  Citações e notas devem ser conforme a ABNT NBR 10520:2002. Segundo Severino (1984, p. 126) ``As citações são referências extraídas de textos durante a pesquisa documental e bibliográfica, consistindo em importantes fontes para a fundamentação das ideias desenvolvidas pelo autor no decorrer de seu raciocínio".

As tabelas devem ser citadas no texto, inseridas o mais próximo possível do trecho a que se referem, e padronizadas conforme as ``Normas de Apresentação Tabular" do Instituto Brasileiro de Geografia e Estatística (IBGE).


\begin{table}[H]
	\centering
	\caption{Título da tabela}
		\begin{tabular}{ccccc}
	\toprule
   \textbf{ORD}	& \textbf{DESCRIÇÃO} & \textbf{XXXX} & \textbf{XXXX} & \textbf{QTD.} \\
	\midrule
	1 & Xxxxx & Abcd & $75,56$ & $3,20$ \\
	2 & Xxxx & Efg & $65,6$ & $34,9585$ \\
	3 & Xxxx & Hij & $34,987$ &  $454,63$ \\
	4 & xxx & lmn & $234,767$ &  $0,982$ \\
	\bottomrule
\end{tabular}
\caption*{Fonte: elaborado pelos autores.}
\end{table}






Deve-se indicar a fonte consultada (elemento obrigatório, mesmo que seja produção do próprio autor), de acordo com a ABNT NBR 10520:2002.





\section{\large RESULTADOS E DISCUSSÕES}


Neste item se faz a descrição detalhada dos resultados obtidos ao longo da pesquisa.  Devem ser apresentadas as informações relevantes obtidas na bibliografia ou nos websites, os resultados dos experimentos realizados em laboratórios, ou seja, todos os dados quantitativos ou qualitativos obtidos ao longo da pesquisa. 









\begin{table}[H]
	\centering
	\caption{Título da tabela}
		\begin{tabular}{ccccc}
		\toprule
		\textbf{ORD}	& \textbf{DESCRIÇÃO} & \textbf{XXXX} & \textbf{XXXX} & \textbf{QTD.} \\
		\midrule
		1 & Xxxxx & Abcd & 75,56 & 1,537 \\
		2 & Xxxx & Efg & 65,6 & 231,33 \\
		3 & Xxxx & Hij & 34,987 & 22,7 \\
		4 & xxx & KLmno & 234,767 & 123,321 \\
		5 & Xxxxx & Abcd & 75,56 & 1,537 \\
		6 & Xxxx & Efg & 65,6 & 231,33 \\
		7 & Xxxx & Hij & 34,987 & 22,7 \\
		8 & xxx & KLmno & 234,767 & 123,321 \\
		9 & Xxxxx & Abcd & 75,56 & 1,537 \\
		10 & Xxxx & Efg & 65,6 & 231,33 \\
		11 & Xxxx & Hij & 34,987 & 22,7 \\
		12 & xxx & KLmno & 234,767 & 123,321 \\
		\bottomrule
	\end{tabular}
\caption*{Fonte: elaborado pelos autores.}
\end{table}





Não se esqueça de sempre dizer em qual fonte foi encontrada cada informação apresentada. Aqui você pode acrescentar algum gráfico e/ou alguma imagem relevante a análise dos resultados.


\begin{figure}[H]
	\centering
	\caption{Título de gráfico}
	\includegraphics[scale=0.6]{fig1}
	\caption*{Fonte: Sobrenome (ano) da fonte consultada.}
\end{figure}



Títulos, fonte e notas de tabelas e ilustrações devem ser inseridos conforme as ``Normas de Apresentação Tabula" do IBGE. 

Os exemplos abaixo ilustram como inserir fórmulas matemáticas. 


\begin{exem} Escrevendo fórmulas matemáticas.

$$\int_0^\infty e^{-st}\,dt=\frac{e^{-st}}{-s}\Bigg|_0^\infty$$
	
\end{exem}


\begin{exem}
	Numerando as fórmulas matemáticas (caso necessário).
	\begin{eqnarray}
	a_2 &=& a_1+r \nonumber \\
	a_3 &=& a_2  +r \Rightarrow a_3= \left(a_1+r\right)+r=a_1+2r \nonumber \\
	a_4 &=& a_3  +r \Rightarrow a_4=\left(a_1+2r\right)+r\Rightarrow a_4 = a_1+3r \nonumber \\
	&\vdots& \nonumber \\
	a_n&=&a_1+\left(n-1\right)r 
\end{eqnarray}	
	
\end{exem}                                                                                                                              

\section{\large CONSIDERAÇÕES FINAIS}

Este é o momento de apresentar as conclusões e os produtos que a pesquisa produziu. É importante apontar quais informações obtidas na pesquisa possibilitou a construção desta conclusão. Também é possível apontar o impacto que a aplicação desta conclusão pode gerar no mundo real. As referências devem seguir as orientações da norma ABNT NBR 6023:2018.




\begin{center}
	\textbf{REFERÊNCIAS}
\end{center}


\begin{flushleft}
A AVENTURA de Darwin. \textbf{Diário do Nordeste}, Fortaleza, 7 jun. 2006. Caderno 3, p. 4.
	
\

ASSOCIAÇÃO BRASILEIRA DE NORMAS TÉCNICAS. \textbf{ABNT NBR 6023}: informação e documentação: referências: elaboração. Rio de Janeiro: ABNT, 2018.  p.74.

\

ASSOCIAÇÃO BRASILEIRA DE NORMAS TÉCNICAS. \textbf{ABNT NBR 10520}: informação e  documentação: citações em documentos: apresentação. Rio de Janeiro: ABNT, 2002. p. 7.


\end{flushleft}



\end{document}
